% ═══════════════════════════════════════════════════════════════
% ARBEITSBLATT: Repaso Unidad 5 + Números 100-1000
% Spanisch Klasse 8 - KOMPAKT (2 A4-Seiten)
% ═══════════════════════════════════════════════════════════════

\documentclass[11pt, a4paper]{article}

\usepackage[utf8]{inputenc}
\usepackage[T1]{fontenc}
\usepackage[ngerman]{babel}
\usepackage{helvet}
\renewcommand{\familydefault}{\sfdefault}

\usepackage[margin=1.5cm, top=1.8cm, bottom=1.5cm]{geometry}
\usepackage{enumitem}
\usepackage{tabularx}
\usepackage{booktabs}
\usepackage{array}
\usepackage{multicol}

\usepackage{xcolor}
\usepackage{tcolorbox}
\tcbuselibrary{skins}

\usepackage{fancyhdr}
\usepackage{lastpage}

% ═══════════════════════════════════════════════════════════════
% FARBEN
% ═══════════════════════════════════════════════════════════════

\definecolor{spanisch}{HTML}{c41e3a}
\definecolor{grau}{HTML}{888888}
\definecolor{hellgrau}{HTML}{f5f5f5}

% ═══════════════════════════════════════════════════════════════
% VARIABLEN
% ═══════════════════════════════════════════════════════════════

\newcommand{\thema}{Repaso Unidad 5 + Números 100-1000}
\newcommand{\fach}{Spanisch}
\newcommand{\klasse}{8}
\newcommand{\maxpunkte}{28}

% ═══════════════════════════════════════════════════════════════
% KOPF- UND FUSSZEILE
% ═══════════════════════════════════════════════════════════════

\pagestyle{fancy}
\fancyhf{}
\renewcommand{\headrulewidth}{0.4pt}
\renewcommand{\footrulewidth}{0pt}

\lhead{\small\fach{} Kl. \klasse}
\chead{\small Arbeitsblatt}
\rhead{\small Seite \thepage/\pageref{LastPage}}

% ═══════════════════════════════════════════════════════════════
% BEFEHLE
% ═══════════════════════════════════════════════════════════════

\newcommand{\punktebox}[1]{\hfill\fbox{\scriptsize \_\kern-1pt/#1}}
\newcommand{\luecke}[1]{\rule{#1}{0.4pt}}
\newcommand{\es}[1]{\textcolor{spanisch}{\textbf{#1}}}

\newtcolorbox{merkkasten}[1][]{
    colback=hellgrau,
    colframe=spanisch,
    boxrule=0.8pt,
    arc=0pt,
    left=6pt, right=6pt, top=5pt, bottom=5pt,
    fonttitle=\small\bfseries,
    title=#1
}

\setlength{\parindent}{0pt}
\setlength{\parskip}{0pt}

% ═══════════════════════════════════════════════════════════════
% DOKUMENT
% ═══════════════════════════════════════════════════════════════

\begin{document}

% ═══════════════════════════════════════════════════════════════
% KOPFBEREICH
% ═══════════════════════════════════════════════════════════════

\begin{center}
    {\large\bfseries\textcolor{spanisch}{Arbeitsblatt: \thema}}
\end{center}

\vspace{3pt}

\begin{tabularx}{\textwidth}{|X|l|l|}
\hline
\textbf{Name:} \luecke{4cm} &
\textbf{Klasse:} \klasse &
\textbf{Datum:} \luecke{2cm} \\
\hline
\end{tabularx}

\vspace{8pt}

% ═══════════════════════════════════════════════════════════════
% AUFGABE 1 + 2 nebeneinander
% ═══════════════════════════════════════════════════════════════

\begin{multicols}{2}

\textbf{\textcolor{spanisch}{Aufgabe 1: La ropa}} \punktebox{5}

\vspace{3pt}
\small
Verbinde mit Linien.

\vspace{5pt}
\begin{tabular}{r@{\hspace{1.8cm}}l}
T-Shirt & el vestido \\[0.8em]
Hose & la camiseta \\[0.8em]
Schuhe & los pantalones \\[0.8em]
Kleid & el jersey \\[0.8em]
Pullover & los zapatos \\
\end{tabular}

\columnbreak

\textbf{\textcolor{spanisch}{Aufgabe 2: Los colores}} \punktebox{4}

\vspace{3pt}
\small
Schreibe die richtige Farbe. Achte auf die Endung!

\vspace{5pt}
\begin{tabular}{ll}
a) la camiseta \luecke{2.5cm} & (weiß) \\[0.6em]
b) el jersey \luecke{2.5cm} & (schwarz) \\[0.6em]
c) los zapatos \luecke{2.5cm} & (braun) \\[0.6em]
d) la falda \luecke{2.5cm} & (rot) \\
\end{tabular}

\end{multicols}

\vspace{6pt}

% ═══════════════════════════════════════════════════════════════
% AUFGABE 3: Modalverben
% ═══════════════════════════════════════════════════════════════

\textbf{\textcolor{spanisch}{Aufgabe 3: Querer, poder, tener que}} \punktebox{5}

\vspace{3pt}
\small
Ergänze die richtige Form des Modalverbs.

\vspace{4pt}
\begin{tabular}{ll@{\hspace{0.8cm}}ll}
a) Yo \luecke{2.2cm} comprar un jersey. & (querer) &
b) Tú no \luecke{2.2cm} ir a la fiesta. & (poder) \\[0.6em]
c) Nosotros \luecke{2.2cm} estudiar. & (tener que) &
d) Ellos \luecke{2.2cm} una chaqueta. & (querer) \\[0.6em]
\multicolumn{4}{l}{e) María \luecke{2.2cm} hablar español. (poder)} \\
\end{tabular}

\vspace{8pt}

% ═══════════════════════════════════════════════════════════════
% MERKKASTEN: Zahlen 100-1000
% ═══════════════════════════════════════════════════════════════

\begin{merkkasten}[Merkkasten: Los números 100-1000]
\small
\begin{tabular}{ll@{\hspace{1cm}}ll@{\hspace{1cm}}ll}
100 & \es{cien} & 400 & \es{cuatrocientos/as} & 700 & \es{setecientos/as} \\
200 & \es{doscientos/as} & 500 & \es{quinientos/as} & 800 & \es{ochocientos/as} \\
300 & \es{trescientos/as} & 600 & \es{seiscientos/as} & 900 & \es{novecientos/as} \\
\multicolumn{2}{l}{} & \multicolumn{2}{l}{} & 1000 & \es{mil} \\
\end{tabular}

\vspace{4pt}
\textbf{Wichtig:} 100 allein = \es{cien} | 100 + Zahl = \es{ciento} (z.B. ciento uno = 101)\\
\textbf{Unregelmäßig:} \es{quinientos} (500), \es{setecientos} (700), \es{novecientos} (900)\\
\textbf{Genus:} Ab 200 maskulin/feminin! (doscientos euros / doscientas personas)
\end{merkkasten}

\vspace{8pt}

% ═══════════════════════════════════════════════════════════════
% AUFGABE 4 + 5 nebeneinander
% ═══════════════════════════════════════════════════════════════

\begin{multicols}{2}

\textbf{\textcolor{spanisch}{Aufgabe 4: Zahlen schreiben}} \punktebox{6}

\vspace{3pt}
\small
Schreibe die Zahlen als spanische Wörter.

\vspace{4pt}
\begin{tabular}{ll}
a) 100 = & \luecke{4cm} \\[0.5em]
b) 350 = & \luecke{4cm} \\[0.5em]
c) 500 = & \luecke{4cm} \\[0.5em]
d) 725 = & \luecke{4cm} \\[0.5em]
e) 999 = & \luecke{4cm} \\[0.5em]
f) 1000 = & \luecke{4cm} \\
\end{tabular}

\columnbreak

\textbf{\textcolor{spanisch}{Aufgabe 5: Zahlen erkennen}} \punktebox{4}

\vspace{3pt}
\small
Schreibe die Zahl als Ziffer.

\vspace{4pt}
\begin{tabular}{ll}
a) doscientos cuarenta & = \luecke{1.5cm} \\[0.7em]
b) quinientos ochenta y tres & = \luecke{1.5cm} \\[0.7em]
c) setecientos quince & = \luecke{1.5cm} \\[0.7em]
d) novecientos noventa y nueve & = \luecke{1.5cm} \\
\end{tabular}

\end{multicols}

\vspace{8pt}

% ═══════════════════════════════════════════════════════════════
% AUFGABE 6: Einkaufsdialog
% ═══════════════════════════════════════════════════════════════

\textbf{\textcolor{spanisch}{Aufgabe 6: De compras -- Einkaufsdialog}} \punktebox{4}

\vspace{3pt}
\small
Ergänze den Dialog auf Spanisch.

\vspace{5pt}

\begin{tabular}{@{}p{2cm}p{13cm}@{}}
\textbf{Vendedor:} & \es{¡Hola! ¿En qué puedo ayudarte?} \\[0.6em]
\textbf{Cliente:} & \luecke{7cm} \hfill \textit{(Ich möchte einen Pullover kaufen.)} \\[0.6em]
\textbf{Vendedor:} & \es{Muy bien. ¿Qué color quieres?} \\[0.6em]
\textbf{Cliente:} & \luecke{7cm} \hfill \textit{(Ich möchte einen blauen.)} \\[0.6em]
\textbf{Vendedor:} & \es{Aquí tienes. El jersey cuesta trescientos euros.} \\[0.6em]
\textbf{Cliente:} & \luecke{7cm} \hfill \textit{(Wieviel kostet das?)} \\[0.6em]
\textbf{Vendedor:} & \es{Trescientos euros.} \\[0.6em]
\textbf{Cliente:} & \luecke{7cm} \hfill \textit{(Gut, ich kaufe ihn!)} \\
\end{tabular}

% ═══════════════════════════════════════════════════════════════
% PUNKTEÜBERSICHT
% ═══════════════════════════════════════════════════════════════

\vspace{6pt}
\hrule
\vspace{4pt}

\begin{flushright}
\textbf{Punkte:} \luecke{1cm} / \maxpunkte
\end{flushright}

\vspace{6pt}

% ═══════════════════════════════════════════════════════════════
% EXTENSION TASKS (★★★)
% ═══════════════════════════════════════════════════════════════

\begin{tcolorbox}[colback=white, colframe=grau, boxrule=0.5pt, arc=0pt,
    left=6pt, right=6pt, top=4pt, bottom=4pt, title={\small\textbf{Extension Tasks (★★★) -- Für Schnelle}}]
\small

\textbf{E1: Rechenaufgabe} \hfill \textit{+2 BE}\\
Pablo kauft eine Jacke für 245€ und Schuhe für 189€.
Schreibe den Gesamtpreis als spanisches Wort.

\vspace{3pt}
\luecke{12cm}

\vspace{6pt}

\textbf{E2: Eigener Dialog} \hfill \textit{+3 BE}\\
Schreibe einen kurzen Einkaufsdialog (4 Zeilen). Du möchtest ein rotes Kleid kaufen, das 599€ kostet.

\vspace{3pt}
\begin{tabular}{@{}p{1.8cm}p{12cm}@{}}
Cliente: & \luecke{11cm} \\[0.5em]
Vendedor: & \luecke{11cm} \\[0.5em]
Cliente: & \luecke{11cm} \\[0.5em]
Vendedor: & \luecke{11cm} \\
\end{tabular}

\end{tcolorbox}

\end{document}
